\section{CausalNex: maintenance, graph semantics and reuse}
\label{exec14:causalnex}

A graph-learning package enters a bank's proposal through two different
questions: whether its mathematical output answers the intended question, and
whether someone can own its continuing operation. CausalNex illustrates why
both questions matter. Its canonical source at revision
\texttt{b87ab2059987}, dated 31 August 2026, declares version 0.12.2. The README
explicitly states that the project reached end of life on 24 June 2026 and
will receive no further updates, bug fixes, security patches or issue and
pull-request review. The later source commit does not rescind that notice.
This is affirmative maintenance evidence from the recovered repository; the
earlier failure to retrieve its historical namespace was only an access
failure.~\citep{exec14_causalnex_primary}

The proposal therefore rejects CausalNex as a default new dependency while
retaining its design as a historical reference. A narrow archival comparison
could still be justified by a specific reproducibility need. It would require
an accountable maintainer, an isolated environment and a declared endpoint;
the old package's existence supplies none of these obligations automatically.
The root notice declares Apache 2.0 and separately excludes rights to the
QuantumBlack trademarks. That notice supports a licensing route, but it is
not a complete inventory of transitive and bundled rights. The packaging
requires Python at least 3.8 and below 3.11, NumPy below 1.24, pandas below 2.0
and pgmpy below 0.1.20, together with other dependencies and Python-specific
scikit-learn requirements. These declarations motivate environment isolation;
they do not establish that a resolved environment is impossible, vulnerable
or inexpensive to maintain.~\citep{exec14_causalnex_primary}

\subsection{A learned structure and an intervention target}

Consider a servicing analyst who records account covariates, an offer and a
subsequent recovery amount. The inspected continuous-data NOTEARS entrypoints
take an $n\times d$ matrix $X$, with accounts in rows and variables in columns.
With $W_{ij}$ representing a coefficient from variable $i$ to variable $j$,
the unpenalized data term and acyclicity function are
\[
 L(W;X)=\frac{1}{2n}\|X(I-W)\|_F^2,
 \qquad h(W)=\operatorname{tr}\{\exp(W\circ W)\}-d.
\]
Here $\circ$ denotes elementwise multiplication, $\exp$ is the matrix
exponential and $\|\cdot\|_F$ is the square root of the sum of squared matrix
entries. The augmented objective adds $\rho h(W)^2/2+\alpha h(W)$, with
penalty $\rho$ and multiplier $\alpha$. The lasso path represents a coefficient
as the difference of nonnegative parts and penalizes their sum. Tabu edges and
tabu parent or child variables become optimizer bounds. These are concrete
source-level computations.~\citep{exec14_causalnex_primary}

Even an acyclic solution supplies a candidate structure under that loss and
those restrictions. An intervention contrast such as expected recovery under
an offer versus no offer additionally needs a causal interpretation, an
eligible population, an outcome horizon and assumptions about confounding and
the observation clock. A post-offer delinquency measurement cannot become a
baseline variable because the optimizer accepts its column. The selected
entrypoints return no causal graph confidence region or treatment-effect
interval. The package also advertises Bayesian-network fitting and
intervention facilities; this selected source review does not qualify those
separate implementations or establish their appropriateness for longitudinal
servicing.~\citep{exec14_causalnex_primary}

The graph class makes another boundary explicit. A \texttt{StructureModel}
can contain directed cycles. Edge origins take the values unknown, learned
or expert; the pandas wrapper preserves column names and marks learned
edges, whereas the selected array-created graph uses the class default.
Origin records who or what supplied an edge. It does not certify its truth.
An assistant can preserve that provenance and ask which temporal or business
restriction supports an expert edge, but should not convert the label
``expert'' into verified causal evidence.~\citep{exec14_causalnex_primary}

\subsection{Preprocessing and graph repair change the reported object}

The selected objective does not center $X$ automatically. To see the
consequence without fitting the upstream package, impose two observations
$(-1,-2)$ and $(1,2)$ for variables $(A,B)$ and restrict the direction to
$A\to B$. With no intercept the least-squares coefficient is
$\sum_i A_iB_i/\sum_i A_i^2=2$, and the full two-variable loss above is $1/2$.
Change the origin of the first variable to $A'=A+1$. The observations become
$(0,-2)$ and $(2,2)$, and their centered second-moment matrix remains
\[
 \begin{pmatrix}1&2\\2&4\end{pmatrix}.
\]
Yet the no-intercept coefficient is now $1$. Holding the old coefficient at
$2$ yields loss $3$; refitting within the fixed direction yields loss $2$.
These are exact calculations for an imposed example, not measurements of
which graph the upstream optimizer would select. They show why a caller must
declare centering, the population used to fit preprocessing and the treatment
of intercepts.~\citep{exec14_causalnex_oracle}

Scaling also affects a subsequent graph operation. The method
\texttt{threshold\_till\_dag} repeatedly deletes the edge with smallest
absolute weight until a DAG remains; its own docstring warns about weights
with different units.~\citep{exec14_causalnex_primary}
For a stipulated two-cycle, let the coefficients be $A\to B=0.2$ and
$B\to A=0.5$. This rule deletes $A\to B$. Express the same coefficient
relations using $B'=10B$: the weights become $2$ and $0.05$, so the rule
deletes $B\to A$. Algebraically, a change of units with diagonal scale matrix
$D$ gives $W'=D^{-1}WD$. In this example the cycle product remains $0.1$ and
both acyclicity values equal $2\cosh(0.1)-2>0$. The cycle has not disappeared
through a change of units; the pruning choice has changed.
\citep{exec14_causalnex_oracle}

In a banking model, two opposite surviving arrows can lead to different
adjustment sets and intervention stories. Deleting a weak edge is thus a
scientific modeling choice, even when it also solves a software input
requirement. A wrapper would need to retain the original graph, the units,
the declared pruning procedure and the full deletion trace. Standardizing
variables may define a more explicit comparison, but cannot prove the
resulting arrows causal or remove observational ambiguity.

Threshold boundaries also deserve an explicit contract. At the inspected
revision the unpenalized path zeros weights satisfying
$|w|\leq t$, while the lasso path zeros weights satisfying $|w|<t$.
For nonzero weights $w=\pm0.2$ and threshold $t=0.2$, the first path removes
both edges and the second retains them. A common adapter that applies a
single threshold convention would silently change one method's output.
Independent rational arithmetic reproduces these boundaries; testing the
actual wrapper would remain a separate implementation task.
\citep{exec14_causalnex_primary,exec14_causalnex_oracle}

\subsection{A conditional archival route and its alternatives}

The inspected optimizer consumes the returned solution vector and can issue a
nonconvergence warning before returning a graph. It does not promote optimizer
success into a structured scientific guarantee. A prospective archival
wrapper would capture termination and warnings, test finite output and
acyclicity, and retain an unqualified candidate when these checks fail. It
would preserve column identities, transforms and tabu directions, and report
timeouts or environment failures as unavailable results. An edge weight would
never be repurposed as a confidence level. These requirements follow from
the selected source and from the meaning of the proposed output; no CausalNex
runtime qualification is claimed here.~\citep{exec14_causalnex_primary}

Direct adoption would transfer unsupported lifecycle obligations to the
product. Porting the optimizer would additionally transfer numerical behavior,
dependency adaptation and testing to the local team; there is no demonstrated
embedding requirement that warrants this ownership. An isolated historical
benchmark is the narrower route when a concrete comparison requires it.
Other discovery engines discussed in this survey need their own pinned
decisions, because constraint tests, score searches and latent-variable
representations answer different questions. Their names are not aliases for
CausalNex, and a separately reviewed pgmpy version does not qualify the older
dependency used here.

For a small banking prototype, an explicit expert DAG and an exactly
enumerated reference model provide useful comparators: they expose assumptions
and intervention arithmetic while leaving discovery accuracy open. The cost
comparison should include retained environments, license and dependency
review, preprocessing, failed requests, analyst correction, future fixes and
retirement. No observed price or maintenance-cost estimate is available from
the present static review. Reconsideration of CausalNex would require a unique
archival need, a named owner, a permissible resolved environment and executed
conversion and failure checks. Public or generated examples can support
those initial checks without purchasing access to bank data.
\label{exec14:causalnex-end}
